\section{Metodología}\label{cap:metodología}

Para el desarrollo de este proyecto se ha decidido hacer uso de una metodología \textbf{iterativa incremental} \cite{sommerville2015}. Este modelo se ha seleccionado porque permite organizar el proyecto en ciclos o iteraciones sucesivas, donde cada una añade funcionalidad al sistema de forma progresiva. Cada iteración incluye sus propias fases de análisis, diseño, implementación y pruebas, lo que facilita la validación continua del producto y la detección temprana de errores.

La elección de esta metodología está justificada por la naturaleza del proyecto, ya que permite entregar incrementos funcionales del sistema de manera gradual, comenzando por las funcionalidades básicas del módulo ciudadano y avanzando progresivamente hacia características más complejas como la geolocalización en tiempo real, los paneles de gestión y el sistema de notificaciones. Al ser un proyecto desarrollado por un único desarrollador, este modelo resulta adecuado porque combina la flexibilidad de adaptarse a posibles ajustes durante el desarrollo con una planificación estructurada que facilita el seguimiento del avance.

El proyecto se ha organizado en una \textbf{fase inicial de análisis} general, seguida de \textbf{cinco iteraciones} que abordan de forma incremental los distintos módulos del sistema, y una \textbf{fase final} de documentación, integración y entrega.